\documentclass[../Main.tex]{subfiles}

\begin{document}
\section{Multivariate Functionals}
We now consider functions:
\begin{align*}
    \vec{y} : V \subseteq \R^m &\mapsto \R^n\\
    \vec{x} &\mapsto \vec{y}(\vec{x})
\end{align*}
and functionals of the form:
\begin{equation*}
    F[\vec{y}] = \int_V f(\vec{y}(\vec{x}), \nabla \vec{y}(\vec{x}), \vec{x}) d^m\vec{x}
\end{equation*}
We generally find the specific version of the Euler-Lagrange equation for each problem, because the general Euler-Lagrange derivation is very complex. The process for this is as follows.
\begin{enumerate}
    \item Calculate the relevant first derivatives of $f$;
    \item use these in the Taylor expansion of $F[\vec{y} + \vec{\delta y}] - F[\vec{y}]$;
    \item split the integral into an exact expression $\nabla \vec{g} \cdot \vec{dx}$ (which becomes a boundary term) and a leftover term;
    \item apply relevant boundary conditions on $\partial V$ to ensure the boundary term vanishes;
    \item evaluate the functional derivative as the leftover term.
\end{enumerate}
\section{Examples}
\subsection{Minimal Surface Equation}
The Minimal Surface Equation gives the function $h(x, y)$ that minimises the area of a surface $z = h(x,y)$ connecting a curve $\partial D$:
\begin{equation*}
    (1 + h_y^2) h_{xx} + (1 + h_x^2) h_{yy} - 2h_x h_y h_{xy} = 0
\end{equation*}
If $h_x, h_y$ are small and positive we derive the Laplace equation. If the boundary curve lines in a plane then we get the planar solution $Ax + By + C = 0$.

If $h = h(r)$ we find the area functional has a first integral, and solving this gives a Catenoid:
\begin{equation*}
    r = r_0 \cos(\frac{z - z_0}{r_0}).
\end{equation*}
\subsection{Vibrating String}
The action of a vibrating string is:
\begin{equation*}
    I[y] = \frac12 \int_{t = t_1}^{t_2} \int_{x=0}^{a} \rho \dot{y}^2 - Ty^{\prime 2} dx~dt 
\end{equation*}
and the Euler-Lagrange equation is the wave equation from IB Methods.
\section{Actions and Maxwell's Equations}
Given the following two Maxwell equations:
\begin{align}
    \nabla \cdot \vec{B} &= 0 \label{eqnMBD} \\
    \nabla \times \vec{E} &= -\frac{\partial \vec{B}}{\partial t} \label{eqnMEC}
\end{align}
then we can derive the other two. Consider potentials $\vec{A}$ and $\phi$ such that:
\begin{equation*}
    B = \nabla \times \vec{A},\qquad \vec{E} + \frac{\partial \vec{A}}{\partial t} = -\nabla \phi
\end{equation*}
and then the action is a functional of $\vec{A}$ and $\phi$:
\begin{equation*}
    I[\vec{A}, \phi] = \int_{t = 0}^T\int_{\vec{x} \in D} \left(\frac{|\vec{E}|^2}{2\epsilon_0} - \frac{\mu_0|\vec{B}|^2}{2} + \vec{A} \cdot \vec{J} - \rho \phi\right) d^3\vec{x} dt
\end{equation*}
Varying this ($\vec{A} \to \vec{A} + \delta\vec{A}$, and the same for $\phi$), integrating by parts and setting the coefficients of $\delta\vec{A}$ and $\delta\phi$ to zero gives the other two Maxwell equations.
\begin{align}
    \nabla \cdot \vec{E} &= \frac{\rho}{\epsilon_0} \label{eqnMED} \\
    \nabla \times \vec{B} &= \mu_0 \vec{J} + \mu_0\epsilon_0 \frac{\partial \vec{E}}{\partial t} \label{eqnMBC}
\end{align}
\end{document}